\scanpage{90}
\margin{89}
\mspage{57}
limit ourselves to the case $g=2$. We are interested in $\mathrm{M}(\Gamma_2(\lambda))$ for $\lambda=1,2$, because we know $\mathrm{S}(\Gamma_2(\lambda))$ explicitly for the same values of $\lambda$ (8). We keep in mind that they are defined as quotient varieties and not as monoidal transforms. We take a point $t_{\unsure{o|0}}$ of $\underline{\mathrm{S}}_1$. Since $t_{\unsure{o|0}}$ is not a fixed point of any element of $\Gamma_1(2)$ different from $\pm1_2$, the analytic local ring $\mathrm{O}$ of $\mathrm{M}(\Gamma_2(2))$ at the image point of $t_{\unsure{o|0}}$ consists of convergent series of the form
\[
\sum_{k=0}^{\infty}\theta_k(t,z)e(kw/2)\ .
\]
The coefficient $\theta_k(t,z)$ is holomorphic in $\mathrm{V}\times\mathbb{C}$ for some $\mathrm{V}$ and $z\longrightarrow\theta_k(t,2z)$ is a theta-function of order $4k$, which is even in the sense it is invariant under $z\longrightarrow-z$. This follows from Theorem 1. Moreover, the ideal $\mathrm{I}$ consists of convergent series which start from $k=1$. Furthermore, there exists a base $(\theta_{\unsure{l|1}}(t_{\unsure{o|0}},z))$ over $\mathbb{C}$ of the vector space of even theta-functions of order four such that each $\theta_{\unsure{l|1}}(t,z)$ is holomorphic in some $\mathrm{V}\times\mathbb{C}$ and such that the correspondence $z\longrightarrow(\theta_{\unsure{l|1}}(t,z))$ gives rise to a complex-analytic embedding of the quotient variety of $\mathrm{T}_1(t)$ by its involution $z\longrightarrow-z$, which is $\mathrm{P}_1(\mathbb{C})$, as a plane curve of order two. Therefore, if we take the monoidal transform of $\mathrm{S}(\Gamma_2(2))$ along its boundary in the neighborhood of the image point of $t_{\unsure{o|0}}$, we get the quotient variety of the monoidal transform of $\mathrm{S}(\Gamma_2(\lambda))$, for any even $\lambda$ at least equal to 3, along its boundary in the neighborhood of the image point of $t_{\unsure{o|0}}$. In the same way, we can conclude that the morphism $\mathrm{M}(\Gamma_2(1))\longrightarrow\mathrm{S}(\Gamma_2(1))$ is a monoidal transformation in the neighborhood of the image point of $t_{\unsure{o|0}}$ provided $t_{\unsure{o|0}}$ is not a fixed point of any element of $\Gamma_1(1)$ different
